\chapter{Conclusiones generales}
\label{ch:conclusiones}

Este trabajo implementó y validó las tareas computacionales \textbf{T1} y
\textbf{T3} del estudio numérico del efecto Mpemba cuántico, cada una en versión
serial y paralela, sobre la cadena de Ising transversa disipativa.

\ideacentral{La tesis (\cref{sec:tesis}) se cumple: el QMpE es geometría espectral
de un generador no normal, así que verificarlo y explotarlo en HPC equivale a
\emph{calcular} su estructura y dinámica de relajación de forma correcta y
escalable. T1 entrega el \textbf{diagnóstico que predice} el efecto (modos lentos,
$a_k$); T3a/T3b entregan la \textbf{simulación} de la relajación en muchos cuerpos.
Los datos confirman las dos exigencias: \emph{correctitud} (todas validadas contra
el oráculo exacto) y \emph{escalabilidad} (la habilitación para explotar).}

\section*{Correctitud}

Los tres métodos se validaron contra un mismo oráculo exacto (diagonalización e
integración densa del Liouvilliano, $N\lesssim4$):
\begin{itemize}
  \item \textbf{T1} (Arnoldi--Lindblad) reproduce los autovalores lentos con error
  $10^{-13}$--$10^{-5}$ y el estado estacionario a $10^{-11}$.
  \item \textbf{T3a} (trayectorias) converge con la ley estadística $M^{-1/2}$
  (pendiente medida $-0.488$) y coincide con la ecuación maestra a $8\times10^{-3}$
  con $M=2\times10^4$.
  \item \textbf{T3b} (TEBD) converge a la solución exacta al aumentar $\chi$,
  alcanzando $6\times10^{-7}$ cuando $\chi$ iguala el rango del sistema.
\end{itemize}

\section*{La lección de Sistemas Distribuidos: la estructura del acoplamiento
determina la escalabilidad}

El resultado más instructivo es el \emph{contraste} entre los tres patrones de
paralelismo, resumido en la \cref{tab:resumen}:

\begin{table}[h]
\centering
\caption{Comparación de los tres patrones de paralelismo (speedup a 4 workers).}
\label{tab:resumen}
\begin{tabular}{lll c}
\toprule
Tarea & Patrón de paralelismo & Tecnología & speedup (4) \\
\midrule
T1  & SpMV \emph{secuencialmente acoplado} & MPI+OpenMP (datos) & $\sim2.8\times$ \\
T3a & trayectorias \emph{independientes}   & MPI+OpenMP (tareas) & $\sim3.7\times$ \\
T3b & compuertas \emph{disjuntas} por capa & hilos (estructurado) & $\sim2.7\times$ \\
\bottomrule
\end{tabular}
\end{table}

\begin{itemize}
  \item \textbf{T1} está limitado porque Arnoldi es secuencialmente acoplado: solo
  el SpMV interno paraleliza, y su naturaleza \emph{memory-bound} satura el
  escalado en el hardware de prueba. La comunicación \texttt{Allgatherv} de MPI se
  amortiza por el cómputo $\mathcal{O}(d^3)$, de modo que MPI iguala a OpenMP.
  \item \textbf{T3a} es \emph{embarrassingly parallel} y escala casi idealmente,
  con OpenMP y MPI indistinguibles; es el mejor candidato a escalar a muchos nodos.
  \item \textbf{T3b} tiene paralelismo estructurado de granularidad media,
  limitado por el número de bonds disjuntos y el desbalance de carga (los tensores
  centrales soportan mayor $\chi$).
\end{itemize}

\section*{Alcance en muchos cuerpos}

Las tareas T3 cumplen su promesa de superar la barrera $4^N$: T3a simula $N=10$
($d=1024$) con trayectorias de memoria $2^N$, y T3b alcanza $N=32$
($4^{32}\approx1.8\times10^{19}$) con dimensión de enlace acotada $\chi=24$. Ambas
son inalcanzables para la representación explícita del Liouvilliano que requiere T1.

\section*{Trabajo futuro}

Las costuras de paralelización pendientes, en línea con el informe base, son:
(i) almacenamiento disperso (CSR) de $H$ y $L_\mu$ para distribuir realmente la
\emph{memoria} en T1 (no solo el cómputo); (ii) trayectorias en GPU para T3a;
(iii) TEBD distribuido por segmentos de cadena con frontera MPI para T3b; y (iv)
combinar T1 y T3b (espectro del Liouvilliano con redes tensoriales) para el gap en
cadenas largas. En conjunto, las tres tareas forman una cadena de herramientas
coherente: T1 diagnostica el efecto en pocos cuerpos, y T3a/T3b lo exploran en el
régimen de muchos cuerpos donde el fenómeno es físicamente más rico.
